\documentclass[11pt,letterpaper]{article}
\usepackage[T1]{fontenc}
\usepackage[utf8]{inputenc}
\usepackage{lmodern,microtype}
\usepackage[margin=1in,headheight=15pt]{geometry}
\usepackage{amsmath,amssymb,amsthm,mathtools}
\usepackage{booktabs,longtable,array,enumitem}
\usepackage{xcolor,graphicx,fancyhdr}
\usepackage[hyphens]{url}
\usepackage{xurl}
\usepackage[colorlinks=true,linkcolor=blue!45!black,citecolor=blue!45!black,urlcolor=blue!45!black]{hyperref}
\usepackage{listings}
\definecolor{ink}{HTML}{17324D}
\definecolor{pale}{HTML}{F1F4F7}
\lstset{basicstyle=\ttfamily\small,breaklines=true,columns=fullflexible,keepspaces=true,
  backgroundcolor=\color{pale},frame=single,rulecolor=\color{black!15},xleftmargin=4pt,xrightmargin=4pt}
\pagestyle{fancy}\fancyhf{}
\fancyhead[L]{\small Matroidal dual-forest selection}
\fancyhead[R]{\small ProveIt research continuation}
\fancyfoot[C]{\thepage}
\setlength{\parskip}{0.35em}
\setlength{\parindent}{0pt}
\setcounter{tocdepth}{2}
\newtheorem{theorem}{Theorem}[section]
\newtheorem{lemma}[theorem]{Lemma}
\newtheorem{proposition}[theorem]{Proposition}
\newtheorem{corollary}[theorem]{Corollary}
\theoremstyle{definition}
\newtheorem{definition}[theorem]{Definition}
\newtheorem{example}[theorem]{Example}
\newtheorem{question}[theorem]{Research question}
\theoremstyle{remark}
\newtheorem{remark}[theorem]{Remark}
\DeclareMathOperator{\supp}{supp}
\DeclareMathOperator{\rank}{rank}
\DeclareMathOperator{\Reach}{Reach}
\DeclareMathOperator{\Anc}{Anc}
\DeclareMathOperator{\cc}{cc}
\DeclareMathOperator{\Bell}{Bell}
\DeclareMathOperator{\poly}{poly}
\newcommand{\N}{\mathbb N}
\newcommand{\F}{\mathbb F}
\newcommand{\code}[1]{\nolinkurl{#1}}
\newcommand{\pin}{\texttt{c6d0ec8aa0b6308310ab6c081ba29d8bc363ede5}}
\hypersetup{pdftitle={Matroidal Dual-Forest Selection for Unknot Recognition},pdfauthor={ChatGPT; prepared for Vladimir Reshetnikov and ProveIt},pdfsubject={Exact local cancellation selection, deletion deficits, fixed-parameter algorithms, and source-replay integration}}
\setlength{\emergencystretch}{2em}
% Generated from results/benchmark.json; see experiments/render_tables.py.
\newcommand{\UniqueLastRatio}{554}
\newcommand{\ObstructedRatio}{381}
\newcommand{\EasyFirstRatio}{0.403}
\newcommand{\LoopRatio}{0.0306}
\newcommand{\LoopSlowdown}{32.7}

\begin{document}
\hypersetup{pageanchor=false}
\begin{titlepage}
\vspace*{0.05in}
{\large\sffamily\color{ink} TOPOLOGY / UNKNOT RECOGNITION\par}
\vspace{0.3in}
{\LARGE\bfseries\color{ink} Matroidal Dual-Forest Selection\par
for Unknot Recognition\par}
\vspace{0.22in}
{\large Polynomial pair-column cancellation, exact deletion deficits,\par
and a fixed-parameter higher-support kernel\par}
\vspace{0.20in}
{Research contribution prepared with ChatGPT\par
for Vladimir Reshetnikov's \textit{ProveIt} project\par}
{12 October 2026\par}
\vspace{0.10in}
\begin{abstract}
The maintained ProveIt normal-cut pipeline can eliminate all source core three-handles by a boundary-rooted dual forest. Its own experiments show that choosing an arbitrary such forest can leave an uncancellable raw core even when a different forest succeeds. We study the exact selection problem, with the original zero/one-handle forest fixed and with unsigned source incidences retained.

For matrices whose usable columns meet at most two rows, deletion closure is directed reachability. A successful dual tree must meet every unanchored source strongly connected component. We reduce that requirement to graphic/partition matroid intersection and derive the exact optimum
\[
 \delta=\max_{J\subseteq[q]}\left(|J|-\rank_D\!\left(\bigcup_{j\in J}C_j\right)\right)
\]
for additional, formally deferred row deletions after the best tree. Both a completing schedule and a rank certificate of the optimum are obtained in polynomial time.

For general matrices, let $b$ be the number of rows incident to usable columns of support at least three. Ordering those rows and grouping shared seed witnesses gives an exact algorithm with at most $b!\Bell(b)$ matroid-intersection subproblems, hence running time $\poly(\Lambda)2^{O(b\log(b+1))}$ for encoded input size $\Lambda$. An explicit satisfiability reduction proves that the unrestricted abstract problem is NP-complete even when every column has support at most three and all entries lie in $\{0,1,2\}$.

The package includes conventional proofs, an independent certificate verifier, 54 passing unit tests, exhaustive and randomized comparisons, a reconstruction of the repository's five-row loop example, and a source-replay integration seam. These are local selection results, not a general quasi-polynomial unknot recognizer. The hardness construction is not asserted to be geometrically realizable, and full native geometry and recognition were not executed in this study.
\end{abstract}
\vfill
{\small\textbf{Status.} Mathematical proofs and executed research kernels; not externally refereed or proof-assistant formalized. Historical priority is not claimed. Two potentially overlapping incoming dual-forest archives were inventoried but their contents were not audited.}\par
\end{titlepage}
\hypersetup{pageanchor=true}
\pagenumbering{roman}
\begingroup
\fontsize{9.5}{10.8}\selectfont
\setlength{\parskip}{0pt}
\tableofcontents
\endgroup
\clearpage
\pagenumbering{arabic}

\section{Research position and the exact contribution}\label{sec:position}
\subsection{The checkpoint being continued}
This article is pinned to ProveIt commit
\begin{center}\small\pin.\end{center}
The immediate inputs are the maintained \code{synthesis/column_cancellation.tex},
\code{synthesis/core_three_handles.tex}, and
\code{synthesis/three_handles_results.tex}, together with the corresponding native producer and source checker \cite{repo-peeling,repo-dual,repo-results,repo-native,repo-checker}.
The relevant distinction is between \emph{eliminating all three-handles} and
\emph{choosing three-handle cancellations that leave a useful lower presentation}.
The first already has a complete sparse source-bound procedure. The second was
explicitly left open in the inspected maintained discussion.

The recorded loop-source experiments are unusually useful: they identify a failure
caused by a choice, rather than by the absence of a certificate. Four possible
first dual edges are eligible; the default edge leaves a stopping set, while each
of the other three permits completion. Enumerating all spanning trees is a valid
finite fallback, but its worst-case size can be exponential. Our goal is to make
this selection step exact without paying for every tree on a structurally broad
class, and to isolate what must be small outside that class.

The contribution is not a new test based on signed exponent sums. The source
matrix records \emph{unsigned geometric intersections}; the value two means at
least two intersections, not zero modulo two. Source topology remains a separate
input obligation. In particular, the same matrix-and-graph optimization may be
applied to many abstract instances which need not arise from an actual normal cut.

\subsection{Relation to the quasi-polynomial target}
The specific target displayed in Lackenby's inspected lecture slides is
$n^{O(\log n)}$, for $n$ crossings \cite{lackenby-talk}. The July 2026 preprint
\emph{Incompressible surfaces, hierarchies and unknot recognition} supplies a
hierarchy framework; its Section~9 distinguishes a bound on iterations from the
still unspecified costs of some operations \cite{lackenby2026}. The present paper
therefore charges its local search explicitly, without inferring a recognition
bound merely from a short hierarchy or a compressed cut.

We prove a polynomial special case, a fixed-parameter generalization, and a
hardness boundary for the abstract selector. Section~\ref{sec:global} states a
conditional composition theorem specifying exactly which further bounds would
make this component compatible with the stated target. It does not assume that
those further bounds are already true.

\subsection{Attribution and novelty boundary}
Stopping-set persistence and complete full-row peeling are maintained results,
recalled in Section~\ref{sec:model}, not new claims. Rooted dual two/three
cancellation is also inherited. The graphical matroid, independent representatives,
and shortest-augmenting-path framework are classical; a modern primary reference
that recalls the relevant intersection machinery is \cite{matroid2024}.

What is developed here is the directed reduction for unsigned pair-column cores,
the resulting exact deletion-deficit interpretation, the exceptional-row ordering
and shared-witness partition reduction, and an explicit support-three hardness
construction for the combined problem. The proofs are given independently of
experimental success. We do not assert first priority in the wider literature.
Moreover, \code{docs/incoming} contained two recently delivered archives with
dual-forest titles whose binary bodies were not inspected. This is a real audit
limitation, not evidence that those archives lack overlapping results. The
package's \code{SOURCE_AUDIT.md} names them and separates their inventory from the
maintained text actually read.

\section{Authenticated cores and the cancellation model}\label{sec:model}
\subsection{Input, conventions, and scope}
Let $E=\{0,\ldots,m-1\}$ be a row set, and let
\[
 N=(n_{rh})_{r\in E,\,h\in H}\in\N^{E\times H}
\]
be a finite nonnegative matrix. Let $D=(V,E)$ be a connected undirected multigraph
with a distinguished vertex $\partial$. A row labels one dual edge, including
possibly a loop. Parallel edges are distinct. Put $v=|V|$ and let $s$ be the number
of nonzero source entries. The input length $\Lambda$ includes all row, column,
and endpoint labels and all coefficient bits. Zero columns can be omitted.

For the intended native application, $V\setminus\{\partial\}$ consists of core
three-handles, and a row is a source two-handle. Original-boundary and cut-boundary
endpoints are identified with $\partial$. The matrix is formed after one fixed
zero/one-handle forest has been contracted. Its columns are the retained
one-handles, and its entries count intersections of attaching curves with their
belt disks. The source must authenticate these facts; arbitrary group relators
are not substitutes for the geometric counts.

\begin{definition}[Dual-forest cancellation selection]\label{def:problem}
A spanning tree $T$ of $D$ is \emph{successful} if $N[E\setminus T,H]$ admits
complete positive one/two cancellation. The selection problem asks whether such
a tree exists and, on success, asks for a tree and a complete schedule.
\end{definition}

A ``no'' answer in Definition~\ref{def:problem} means only that this fixed
presentation and move family cannot finish. It is not a certificate that a
manifold is not a handlebody, or that a knot is nontrivial. A different primal
forest, geometric intersection reduction, or new cutting surface may change the
answer.

\subsection{The geometric interface inherited from the repository}
\begin{proposition}[Source-bound dual-tree interface; inherited]\label{prop:geometry}
Assume the authenticated core charts satisfy the endpoint description in
\cite{repo-dual}. Every spanning tree of $D$ can be canceled in root-to-leaf order
by geometric two/three-handle cancellations. The surviving two-handle attaching
maps in the original zero/one body, and the bundle port charts in that body, are
unchanged. Thus its effect on the fixed raw matrix is precisely deletion of the
tree rows.
\end{proposition}
\begin{proof}
For a tree edge from the already exposed root component to a new vertex, its two
endpoint occurrences include exactly one live three-handle. It therefore supplies
the authenticated single-intersection cancellation. Exposing the new vertex makes
it part of the root component for subsequent steps. Repeating along the tree
removes all $v-1$ interior vertices. A loop has two equal live occurrences or no
live occurrence, and never supplies a tree edge. Preservation of the lower body
and ports is the geometric relative-cancellation statement established in the
maintained source argument \cite{repo-dual}; it is an assumption about those
charts, not a conclusion from an arbitrary matrix.
\end{proof}

The source endpoint rule is especially transparent on an original edge with
normal weight $w$. For $w=0$, its gap has the two original vertex endpoints. For
$w>0$, the two extreme gaps each join one original endpoint to the cut boundary.
This accounts for both repeated endpoints and edges directly to the boundary.
The native independent replayer reconstructs this information from link triangles
and local normal coordinates, rather than trusting the producer's graph
\cite{repo-checker}.

\subsection{Why full positive cancellation is a peeling question}
A positive pivot at $(r,g)$ has $n_{rg}=1$. After deleting its row and column, the
other unsigned incidences obey
\begin{equation}\label{eq:positive-update}
 n'_{ih}=n_{ih}+n_{ig}n_{rh}.
\end{equation}
A \emph{unit-column peel} is the case where column $g$ has total one on all live
rows. It creates no fill in surviving rows.

\begin{definition}[Stopping set]
A nonempty $R\subseteq E$ is a stopping set if, for every column $h$,
\[
 S_h(R):=\sum_{r\in R}n_{rh}\quad\text{is either zero or at least two.}
\]
\end{definition}

\begin{lemma}[Stopping persistence; inherited]\label{lem:stopping}
A nonempty stopping set prevents full positive cancellation, even if nonzero-fill
pivots are allowed. Maximal unit-column peeling deletes all rows if and only if
full positive cancellation is possible. Its residual is the unique largest source
stopping set.
\end{lemma}
\begin{proof}
We reproduce the maintained argument \cite{repo-peeling} to identify the exact
move system. If a pivot row $r$ is outside $R$, then for a surviving column $h$,
\[
 S'_h(R)=S_h(R)+S_g(R)n_{rh}.
\]
The right side is zero or at least two. If $r\in R$, then $S_g(R)\ge2$, so
$R'=R\setminus\{r\}$ is nonempty, and
\[
 S'_h(R')=S_h(R)+(S_g(R)-2)n_{rh}.
\]
When $S_h(R)=0$, also $n_{rh}=0$; otherwise the displayed sum is at least two.
A nonempty stopping set therefore survives every positive-pivot prefix.

Peeling creates no fill. If it stalls on $R$, every surviving column has total
zero or at least two on $R$. Every previously deleted column is zero on $R$,
because it met only its deleted pivot row when used. Thus $R$ is a source stopping
set. Conversely, a peel cannot remove the first row of any source stopping set.
Finally, a union of stopping sets is a stopping set: a column meeting the union
already meets a constituent with total at least two. All source stopping sets lie
in the residual, proving maximality and uniqueness.
\end{proof}

\begin{lemma}[Monotone seed deletion and saturation]\label{lem:monotone}
If deleting $F\subseteq E$ permits full peeling, deleting any $F'\supseteq F$
also does. Replacing each coefficient by $\min(2,n_{rh})$ preserves every peeling
answer and, consequently, every full positive-cancellation answer.
\end{lemma}
\begin{proof}
A nonempty stopping set contained in $E\setminus F'$ would also be a source
stopping set contained in $E\setminus F$, contradicting Lemma~\ref{lem:stopping}.
For saturation, the predicate that a nonnegative column sum is one depends only
on which entries are zero, which are one, and which exceed one. Every successive
static peel is therefore unchanged. Apply Lemma~\ref{lem:stopping} again.
\end{proof}

It follows that a successful \emph{independent} dual edge set suffices: extend it
to a spanning tree, and the extra deletions cannot spoil success. This elementary
observation is what connects cancellation closure to a matroid constraint.

\section{Pair columns become directed reachability}\label{sec:reachability}
\subsection{Usable columns and the dependency digraph}
A column is \emph{usable} if at least one of its entries is one. A column with no
unit entry can never be a static peeling pivot, regardless of how many other rows
are deleted. It may therefore be omitted from the closure calculation.

Initially suppose every usable column has support size at most two. Construct a
digraph $G_N$ on the row set $E$ and an anchor set $A$ as follows. A singleton
usable column on row $r$ places $r$ in $A$. For a column supported on
$\{u,w\}$, put an arc $w\to u$ when $n_{uh}=1$, and put an arc $u\to w$ when
$n_{wh}=1$. A pair with two unit entries gives two opposite arcs. A pair with
entries one and two gives just the arc from the double-incidence row to the
unit-incidence row.

For $S\subseteq E$, write $\Reach(S)$ for reachability in $G_N$, including $S$
itself. Let $\mathcal P(F)$ denote all rows deleted by first deleting $F$ and
then peeling to exhaustion, including the initial rows $F$.

\begin{theorem}[Exact directed closure]\label{thm:reach}
For a pair-column input and every initial deletion set $F$,
\begin{equation}\label{eq:reach}
 \mathcal P(F)=\Reach(A\cup F).
\end{equation}
\end{theorem}
\begin{proof}
A usable singleton column removes its anchored row unless that row is already
gone. A usable pair column can remove a row $u$ exactly when its other row has
already been removed and its entry at $u$ is one. These are precisely the anchor
and directed-arc rules. Thus each actual peel stays in the right side of
\eqref{eq:reach}. Conversely, follow any directed path starting at an anchor or
an initially deleted row. Once the predecessor is deleted, the column supplying
the next arc removes its unit row if that row is still present. Induction on path
length proves the reverse inclusion.

Deleting a pivot column does not invalidate this reasoning. When it is used, all
its other incident rows are already gone; it has no remaining target whose future
deletion might depend on reusing it.
\end{proof}

\subsection{Only unanchored source components require seeds}
Contract the strongly connected components of $G_N$ to obtain its condensation
DAG. A \emph{source component} has no incoming arc from another component. Let
\[
 C_1,\ldots,C_q
\]
be the source components containing no anchor. These row sets are disjoint.

\begin{corollary}[Necessary and sufficient seed blocks]\label{cor:blocks}
For a pair-column input, deleting $F$ permits full cancellation if and only if
\begin{equation}\label{eq:block-hit}
 F\cap C_i\ne\varnothing\qquad(1\le i\le q).
\end{equation}
\end{corollary}
\begin{proof}
An unanchored source component cannot be reached from another component, so it
needs an initial seed in itself. Conversely, every component in a finite DAG is
reachable from a source. A source containing an anchor is already reached; a
source meeting $F$ is reached as well. Equation~\eqref{eq:reach} completes the
argument.
\end{proof}

\begin{example}[The direction matters]
For the one-column matrix $(1,2)^{\mathsf T}$, the only arc is $1\to0$.
Deleting row~1 clears the matrix, while deleting row~0 leaves the entry two.
Replacing this digraph by an undirected support graph would incorrectly make
both seeds equivalent. An anchor farther downstream also does not seed an
unanchored source: an arc $0\to1$ and an anchor at~1 still require a seed at~0.
\end{example}

\begin{remark}[This is support size, not matrix rank]
The condition concerns the number of rows touched by each usable column. It is
unrelated to the rank of $N$ over a field. In particular, no reduction of its
unsigned coefficients modulo two is performed.
\end{remark}

\section{Independent representatives with checkable rank obstructions}\label{sec:intersection}
\subsection{The common combinatorial primitive}
We need a primitive slightly more general than the disjoint blocks above.
Let $Q_1,\ldots,Q_k\subseteq E$ be arbitrary candidate sets. A partial independent
transversal consists of distinct representatives, one for each of some colors,
whose dual edges form a forest. Let $\nu(Q)$ be the maximum number of represented
colors. For $X\subseteq E$, the graphic rank is
\begin{equation}\label{eq:graphic-rank}
 r_D(X)=v-\cc(V,X).
\end{equation}
Loops contribute no rank; parallel edges cannot both belong to an independent set.

Introduce one \emph{clone} $(r,i)$ for each $r\in Q_i$. Clones of a row share its
dual endpoints and are parallel in the graphic matroid. A partition matroid
permits at most one clone of each color $i$. Consequently a common independent
set is exactly a partial independent transversal. The number of clones is
$P=\sum_i|Q_i|$.

\begin{theorem}[Independent representatives and exact deficiency]\label{thm:rado}
There is a polynomial algorithm returning a maximum partial independent
transversal and a color subset $J\subseteq[k]$ such that
\begin{equation}\label{eq:rado-deficiency}
 k-\nu(Q)=|J|-r_D\!\left(\bigcup_{i\in J}Q_i\right)
 =\max_{K\subseteq[k]}\left(|K|-r_D\!\left(\bigcup_{i\in K}Q_i\right)\right).
\end{equation}
In particular, a full independent transversal exists if and only if every color
subset $K$ satisfies the corresponding rank inequality.
\end{theorem}
This is the classical independent-representative/matroid-intersection conclusion.
We give its specialized algorithm and certificate proof so that the delivered
checker does not need to trust an optimizer.

\subsection{Shortest augmentations in the graphic/partition instance}
Maintain a selected clone set $I$ which is both a forest and color-distinct. An
unselected clone $x$ is a source if it can be added to the forest, and a sink if
its color is unused. For selected $y$ and unselected $x$, put an arc $y\to x$
when $I-y+x$ is a forest. Put an arc $x\to y$ when $y$ occupies the color of $x$.
For a source, incoming arcs can be omitted; for a sink, outgoing arcs can be
omitted. Neither omission changes a shortest source-to-sink path.

If the endpoints of $x$ are already connected in $I$, the selected edges $y$
allowing $y\to x$ are exactly the edges on their unique forest path. A loop has
an empty path and is never addable. Thus all exchanges are found with ordinary
forest traversals. A shortest source-to-sink path is found by breadth-first
search. Toggle its selected and unselected vertices.

\begin{lemma}[A shortest toggle is an augmentation]\label{lem:augment}
The toggle of a shortest source-to-sink path preserves the forest and color
conditions and increases $|I|$ by one.
\end{lemma}
\begin{proof}
Write the path as
\[
 x_0,y_1,x_1,\ldots,y_t,x_t,
\]
where the $x_i$ are unselected. For $i<t$, the color of $x_i$ is the color of
$y_{i+1}$; the terminal color of $x_t$ is unused. All these occupied colors are
distinct because the selected vertices of a shortest path are distinct. Deleting
the $y_i$ and inserting the $x_i$ therefore preserves the color condition.

For the graphic condition, $x_0$ joins two different components of the original
forest. Each $x_i$ with $i\ge1$ joins vertices in the same original component;
otherwise it would itself be a source and shorten the path. Its fundamental
cycle in $I+x_i$ contains $y_i$. It contains no $y_j$ with $j<i$, since the arc
$y_j\to x_i$ would shorten the chosen path. The matrix of these fundamental
cycles on the deleted edges $(y_1,\ldots,y_t)$ is triangular over $\F_2$ with
unit diagonal. It is therefore nonsingular.

All cycles in $I+\{x_0,\ldots,x_t\}$ are combinations of the fundamental cycles
of $x_1,\ldots,x_t$: $x_0$ is the only added edge joining its original components
and is not in a cycle. A nonzero combination cannot be zero on all the $y_i$ by
the triangularity just proved. Hence after those edges are deleted no cycle
survives. The net increase in cardinality is one. The use of $\F_2$ here concerns
graph cycle space only, not the source incidence matrix $N$.
\end{proof}

\subsection{Why an unsuccessful search proves optimality}
\begin{proof}[Proof of Theorem~\ref{thm:rado}]
Repeat Lemma~\ref{lem:augment} until there is no augmenting path. Let $R$ be the
set of clone vertices reached from all sources in the last exchange search.
Every reached unselected clone has an occupied color whose selected clone is
also reached, since no sink is reached. Thus the number of colors appearing in
$R$ is $|I\cap R|$.

Every unselected clone outside $R$ is spanned by $I\setminus R$: it is not a
source, and its fundamental forest path contains no selected edge in $R$, or the
exchange arc would make it reached. Therefore the graphic rank of all clones
outside $R$ is $|I\setminus R|$.

Take $J$ to be the colors not appearing in $R$. All clones of these colors are
outside $R$, and hence
\[
 |J|-r_D\!\left(\bigcup_{i\in J}Q_i\right)
 \ge k-|I\cap R|-|I\setminus R|=k-|I|.
\]
For any color subset $K$, the selected transversal contains at least
$|I|-(k-|K|)$ independent representatives in its union. Consequently
\[
 |K|-r_D\!\left(\bigcup_{i\in K}Q_i\right)\le k-|I|.
\]
Equality holds for $J$. The same counting argument applied to any proposed
transversal bounds its size by $k-|J|+r_D(\bigcup_{i\in J}Q_i)=|I|$.
This proves optimality and the displayed min--max equality.
\end{proof}

The returned certificate needs only the colors $J$. A checker forms their union
and counts graph components to verify the rank. It does not reproduce the
augmenting paths or trust a matrix-rank computation. A positive partial
transversal is checked simply for membership, distinct colors, and acyclicity.

\begin{proposition}[A conservative explicit algorithmic bound]\label{prop:mi-cost}
After dense index normalization, the above forest-path implementation uses
$O(P(v+1)^2+k+v)$ structural work, apart from input preparation, and polynomial
storage. All operations involve labels or counts of polynomial bit length.
\end{proposition}
\begin{proof}
There are at most $v-1$ augmentations. In a search, each of $P$ clones requires at
most an $O(v)$ forest traversal, and the resulting exchange graph has $O(Pv)$
arcs. Its breadth-first search takes the same order of work. Summing over at most
$v$ searches gives the bound. The degenerate graph $v=1$ has no augmentations;
its clones are all loops. Empty color classes are retained when forming the
obstruction. Standard direct arrays after relabeling give the structural model;
the supplied readable Python implementation uses dictionaries and sets, so
individual hash-lookup timings are not being asserted as adversarial constant
bit costs. The resulting worst-case bit cost is still polynomial.
\end{proof}

\section{Polynomial forest choice and the exact deletion deficit}\label{sec:deficit}
\begin{theorem}[Polynomial pair-column selector]\label{thm:pair-selector}
Suppose every usable column of $N$ has support at most two. Let $C_1,\ldots,C_q$
be the unanchored source components of $G_N$. A successful dual tree exists if
and only if
\begin{equation}\label{eq:forest-rado}
 r_D\!\left(\bigcup_{i\in J} C_i\right)\ge |J|
 \qquad\text{for every }J\subseteq[q].
\end{equation}
This decision and a completing tree and peeling schedule can be computed in
polynomial time. A failed instance has a polynomial-size rank obstruction.
\end{theorem}
\begin{proof}
A successful tree meets each $C_i$ by Corollary~\ref{cor:blocks}. Choose one of
its edges in each block. The representatives are distinct because the blocks
are disjoint, and they form a forest as a subset of a tree. This proves the
necessary independent-transversal condition.

Conversely, a full independent transversal extends to a spanning tree of the
connected graph $D$. That tree meets every $C_i$ and hence succeeds. Apply
Theorem~\ref{thm:rado} to decide and construct the transversal. Generate the
actual schedule by peeling the full source matrix after deleting the resulting
tree. If the answer is negative, the returned $J$ violates
\eqref{eq:forest-rado}.
\end{proof}

The criterion is stronger than the cardinality test $q\le v-1$. Different seed
blocks may compete for the same graphical rank. For example, two singleton
blocks represented by parallel dual edges have total rank one, so no tree can
hit both even though each edge is individually eligible.

\subsection{An optimum, not only a yes/no test}
\begin{definition}[Additional deletion deficit]
For a spanning tree $T$, permit an additional arbitrary set $Z\subseteq E\setminus T$
of row deletions before peeling, and put
\[
 \delta_D(N)=\min_{T\text{ tree}}\min\{|Z|:\mathcal P(T\cup Z)=E\}.
\]
These additional deletions are a \emph{formal measure of missing progress}. They
are not automatically legal handle moves. A geometric algorithm wishing to use
them must separately justify what deletes or changes those rows.
\end{definition}

\begin{theorem}[Exact pair-column deficit]\label{thm:deficit}
For a pair-column input,
\begin{equation}\label{eq:deficit}
 \boxed{\displaystyle
 \delta_D(N)=q-\nu(C)
 =\max_{J\subseteq[q]}\left(|J|-r_D\!\left(\bigcup_{i\in J}C_i\right)\right).}
\end{equation}
One can output an optimal tree $T$, an optimal extra-deletion set $Z$, a complete
peeling schedule after $T\cup Z$, and a color subset $J$ proving the lower bound,
all in polynomial time.
\end{theorem}
\begin{proof}
For a fixed tree, every unhit source block needs a further deletion within itself;
no deletion in a different component reaches it. Conversely, one row in each
unhit block suffices. Thus its additional-deletion cost is exactly the number of
blocks it misses.

The maximum number of blocks a tree can hit is $\nu(C)$. Any tree hitting $h$
blocks contains an independent partial transversal of size $h$. Conversely,
a maximum partial transversal extends to a tree hitting at least $\nu(C)$ blocks;
it cannot hit another block without containing a larger independent partial
transversal. This proves $\delta=q-\nu(C)$. Equation~\eqref{eq:rado-deficiency}
gives the rank formula. Choose one extra row from each unhit block, peel the
source matrix, and retain the final rank witness.
\end{proof}

A certificate consists of the tree, the extra rows, the pivots, and a color
subset $J$ for which $|J|-r_D(\bigcup_{i\in J}C_i)=|Z|$. The first three objects
prove the upper bound; the rank equality proves the matching lower bound. In
particular, a positive deficit has a meaningful interpretation even when the
core is geometrically a handlebody: it measures the inadequacy of this raw
presentation and these tree deletions, not the topology.

\begin{example}[Zero rows and dual loops]
With two parallel dual edges and two zero matrix rows, a tree deletes only one
row. The other is a stopping set, so $\delta=1$. If an unanchored source block
consists entirely of dual loops, its rank is zero and it forces positive deficit.
The code retains both zero rows and loops; silently discarding either would
invalidate the criterion.
\end{example}

\section{A fixed-parameter extension for higher-support columns}\label{sec:fpt}
\subsection{The exceptional row boundary}
Return to arbitrary $N$. Build $G_N$ and $A$ using only usable columns of support
one or two. Let $\mathcal H$ be the remaining usable columns, and define
\begin{equation}\label{eq:boundary}
 B=\bigcup_{h\in\mathcal H}\supp(h),\qquad b=|B|.
\end{equation}
The union includes \emph{all} rows of such a column, including its entries equal
to two. A high-support column with no unit entry is not in $\mathcal H$ and does
not contribute to $B$.

Call a source strongly connected component of $G_N$ \emph{mandatory} if it
contains neither an anchor nor a row of $B$. Write these components as
$U_1,\ldots,U_a$, and put
\begin{equation}\label{eq:base-reach}
 R_0=\Reach\!\left(A\cup\bigcup_{j=1}^a U_j\right).
\end{equation}
A mandatory component has no incoming low-column trigger and no high-column
pivot at any of its rows. Every successful initial deletion set must therefore
meet it. Once it has one seed, its whole component and all its low-column
descendants are reached. This explains the symbolic use of the full $U_j$ in
\eqref{eq:base-reach}; those rows are not being granted free unproved deletions.

For a row $x$, write
\[
 \Anc(x)=\{r\in E:x\in\Reach(\{r\})\}.
\]
For each ordering $\pi=(x_1,\ldots,x_b)$ of $B$, process its prefixes
$P_i=\{x_1,\ldots,x_{i-1}\}$. Declare $x_i$ \emph{automatic} at that prefix if
either
\begin{equation}\label{eq:auto-low}
 x_i\in R_0\cup\Reach(P_i),
\end{equation}
or there exists $h\in\mathcal H$ with
\begin{equation}\label{eq:auto-high}
 n_{x_i h}=1,\qquad \supp(h)\setminus\{x_i\}\subseteq P_i.
\end{equation}
For every nonautomatic row, impose a special requirement
\begin{equation}\label{eq:ancestor-demand}
 F\cap\Anc(x_i)\ne\varnothing.
\end{equation}
There are at most $b$ special requirements.

\begin{lemma}[Separation of mandatory and special requirements]\label{lem:disjoint}
Every special candidate set $\Anc(x_i)$ is disjoint from $\bigcup_jU_j$.
\end{lemma}
\begin{proof}
If a row of some $U_j$ were an ancestor of $x_i$, then $x_i\in R_0$ and
\eqref{eq:auto-low} would make it automatic.
\end{proof}

This disjointness is essential. It prevents one edge from having to count both
as a mandatory representative and as a shared special witness in the reduction
below.

\subsection{Why ordering only exceptional rows is complete}
\begin{theorem}[Exceptional-order characterization]\label{thm:orders}
For an arbitrary matrix, a successful dual tree exists if and only if there is
an ordering of $B$ and a graphic-independent set $F$ that meets every mandatory
$U_j$ and every special candidate set produced by that ordering.
\end{theorem}
\begin{proof}[Sufficiency]
Extend $F$ to a tree if necessary. By monotonicity it is enough to reason with
$F$. Its intersections with mandatory components, together with anchors, make
all of $R_0$ available through low-column closure.

Induct on the order of $B$. If $x_i$ satisfies \eqref{eq:auto-low}, it follows
from anchors, mandatory seeds, or earlier exceptional rows. If it satisfies
\eqref{eq:auto-high}, all other rows of that column are already deleted and its
unit entry supplies a valid pivot unless $x_i$ is already gone. Otherwise a seed
in $F\cap\Anc(x_i)$ reaches it through low columns. Consequently every row of
$B$ is eventually deleted.

Every source component of $G_N$ is now reached: it contains an anchor, is
mandatory and seeded, or contains a row of $B$. Every component lies downstream
of a source. Low-column closure therefore clears all remaining rows. As in
Theorem~\ref{thm:reach}, a column used earlier cannot still be needed to delete a
new row, because all its other rows were already gone when it was used.
\end{proof}
\begin{proof}[Necessity]
Suppose a tree $T$ succeeds, and choose an actual complete peeling schedule after
its rows are deleted. List the rows of $B\cap T$ first in arbitrary order, then
the other rows of $B$ in their deletion order. The resulting list is $\pi$.
Every mandatory component must meet $T$, for the reason given above.

Consider a nonautomatic $x_i$. If it belongs to $T$, then $x_i\in\Anc(x_i)$
already supplies the required witness. If it was removed by a high-column pivot,
all the other rows of that column lie in the earlier exceptional prefix, making
\eqref{eq:auto-high} true, a contradiction. In the remaining case it was removed
through low-column closure. Trace that closure backwards to an anchor, an
initially deleted row, or a previously high-pivoted row. An anchor or a mandatory
seed makes \eqref{eq:auto-low} true, as does an earlier high-pivoted row because
that row lies in $P_i$. Therefore a nonautomatic $x_i$ must have an ancestor in
$T$. The tree itself is the required independent set $F$.
\end{proof}

The order is a sufficient-availability certificate, not a demand for exact
physical timestamps. A seed may make a nominally later row available early.
That can only help sufficiency; necessity uses the actual deletion order, so
conservative prefix tests do not lose all valid orders.

\subsection{Grouping witnesses rather than guessing their row names}
For a fixed ordering, denote the special candidate sets by
$Q_1,\ldots,Q_k$, with $k\le b$. A single selected row may lie in several of
these sets and satisfy all of their requirements. Requiring distinct
representatives for every special set would be too strong.

\begin{lemma}[Shared-witness partition reduction]\label{lem:partition}
Assume the disjoint mandatory sets $U_1,\ldots,U_a$ are disjoint from every
$Q_i$. There is a graphic-independent set hitting all mandatory and special
sets if and only if some set partition $\Pi$ of $[k]$ has a full independent
transversal for the family
\begin{equation}\label{eq:partition-groups}
 U_1,\ldots,U_a,\quad
 Q_K:=\bigcap_{i\in K}Q_i\quad(K\in\Pi).
\end{equation}
\end{lemma}
\begin{proof}
For a feasible independent set, choose one witness for each special requirement
and group the indices which chose the same row. Distinct groups then have
distinct representatives in their intersections. Choose a representative in
each mandatory block. Mandatory representatives are distinct from each other
and from every special witness, by the disjointness assumptions. All these
representatives belong to the original independent set.

Conversely, a full independent transversal of \eqref{eq:partition-groups} meets
every mandatory block. A representative in $Q_K$ simultaneously meets every
special requirement indexed by $K$. Their union is the desired independent set.
\end{proof}

For example, two requirements with the same singleton candidate set $\{e\}$
need only one independent edge. The partition with one block captures this;
forcing two differently colored clones of $e$ would incorrectly demand a
parallel-edge cycle. This is why simply invoking ordinary matroid intersection
on all ancestor sets is not a correct general algorithm.

\subsection{The algorithm and its bit bound}
\begin{theorem}[Fixed-parameter exact forest selection]\label{thm:fpt}
Definition~\ref{def:problem} can be decided in
\begin{equation}\label{eq:fpt}
 \poly(\Lambda)\,b!\Bell(b)
 =\poly(\Lambda)\,2^{O(b\log(b+1))}
\end{equation}
bit operations, where $b$ is defined in \eqref{eq:boundary}. A positive answer
has a tree-and-pivot certificate of polynomial size. A negative answer can be
certified by at most $b!\Bell(b)$ independently checkable rank obstructions,
with the corresponding branches regenerated from the source.
\end{theorem}
\begin{proof}
Construct the low digraph, its source components, and the exceptional boundary.
For each of the $b!$ orders, generate its at most $b$ special candidate sets.
Enumerate their set partitions. For each partition, solve the independent-
transversal instance in \eqref{eq:partition-groups} using
Theorem~\ref{thm:rado}. Theorem~\ref{thm:orders} and
Lemma~\ref{lem:partition} prove that a positive branch exists exactly when the
original selection problem is positive. Extend its representatives to a tree
and peel the full original matrix to produce the final certificate.

Mandatory blocks are disjoint, so they contribute at most $m$ clones. There are
at most $b$ special partition blocks, each contributing at most $m$ clones.
Thus every branch has $P\le m(b+1)$, and its intersection cost is polynomial
by Proposition~\ref{prop:mi-cost}. Reachability, ancestor computation, and set
intersection also have polynomial cost. Since $\Bell(k)\le\Bell(b)$ for
$k\le b$, the number of branches is at most $b!\Bell(b)$. For $b\ge1$, both
$b!$ and $\Bell(b)$ are at most $b^b$, proving the exponential form; at $b=0$
there is one branch.

Coefficient bits are read once and coefficients are saturated by
Lemma~\ref{lem:monotone}. No arithmetic on expanded normal-sheet multiplicities
is required by this finite core problem. Label manipulation adds polynomial bit
factors. A failed branch returns a color subset whose rank is too small.
Regenerating every order and every canonical partition and checking the supplied
obstruction for each proves a negative answer without running an optimizer.
\end{proof}

Orders and partitions can be generated depth-first with polynomial workspace.
The implementation can suppress the negative trace with
\code{retain_negative=False}; the default stores it, and may therefore use
$\poly(\Lambda)b!\Bell(b)$ output storage. A streaming proof writer would retain
polynomial working memory while writing a potentially large negative proof.
No polynomial-size negative certificate for all instances is claimed.

The parameter is the size of a higher-support \emph{row boundary}, not the
number of higher-support columns. One large usable column may make $b$ large.
Nor is merely bounding each column by three enough to make $b$ small, as the next
section shows.

\section{A sharp abstract hardness boundary}\label{sec:hardness}
\begin{theorem}[NP-completeness at support three]\label{thm:hardness}
The abstract problem in Definition~\ref{def:problem} is NP-complete even under
all of the following restrictions: every nonzero entry is one or two; every
column has support at most three; and the dual graph is a rooted star with two
parallel edges on each arm, together with loops at the root. The hard instances
can have no zero rows and no singleton columns.
\end{theorem}
\begin{proof}
Membership in NP follows from a spanning tree and a complete timestamp-checked
peeling schedule. For hardness, reduce from Boolean satisfiability in conjunctive
normal form, whose NP-completeness is the classical Cook--Levin foundation
\cite{cook1971}. We may use instances with at least two nonempty clauses; harmless
duplicated clauses give this restriction without affecting hardness.

For variable $x_i$, create two row edges $t_i,f_i$ joining the root to a new
vertex $v_i$. Every spanning tree chooses exactly one of them. Interpret these
choices as the truth values. All other rows constructed below are root loops,
so they can never be initial tree deletions.

An implication from previously deleted rows $a_1,\ldots,a_d$ to a target $z$
is encoded by a column with entry one at $z$, entries two at the $a_j$, and zeros
elsewhere. The target becomes peelable exactly after all prerequisites are
deleted. No prerequisite can itself pivot at this column.

Create a loop row $c_j$ for each clause. For every literal in that clause, add a
one-prerequisite column from its corresponding literal row ($t_i$ or $f_i$) to
$c_j$. These columns implement the disjunction of the chosen literals. Combine
all clause rows by a binary tree of two-prerequisite implication columns,
creating new loop rows for its internal AND gates. Let $z$ be its root, the
accept row. Finally, add a one-prerequisite implication from $z$ to every literal
row. These cleanup columns allow the unchosen literal rows to disappear, but
only after acceptance.

Every column has support two or three and entries in $\{0,1,2\}$. With nonempty
clauses every clause row has a unit entry; every literal row has a unit cleanup
entry; every AND row has a unit defining entry. Hence there are no zero rows.
The number of rows and columns is polynomial in the formula length.

If a satisfying assignment is given, delete its selected literal tree rows.
Every clause is then triggered, the AND tree evaluates to its accept row, and
the cleanup columns remove all unchosen literals. All rows are deleted.

Conversely, consider the time at which the accept row first disappears in a
successful peeling schedule. Before that time, an unchosen literal cannot have
disappeared: its only columns in which it has entry one are cleanup columns
requiring acceptance. Each deleted clause must therefore have had a literal
chosen by the tree. Acceptance requires all leaves of the AND tree to have
been deleted, so all clauses are satisfied by the tree's assignment. Thus the
formula is satisfiable exactly when its matrix instance has a successful tree.
Lemma~\ref{lem:stopping} extends this equivalence from peeling to the permitted
full positive-cancellation rule.
\end{proof}

\paragraph{The geometric limitation is substantive.}
Theorem~\ref{thm:hardness} is not a hardness result for unknot recognition, for
geometric handlebody recognition, or even for the class of source-derived raw
cores in the repository. We have not realized these matrices and dual graphs by
normal-cut attaching charts. Geometric compatibility may exclude the construction
or force additional structure. What the theorem does show is that \emph{the
abstract incidence and dual-graph conditions alone} do not justify a polynomial
selector for all bounded-arity inputs unless P equals NP.

The reduction also explains the FPT parameter: usable three-row columns form
binary AND gates, while low columns transmit information through arbitrarily
large directed networks. The exceptional boundary isolates where genuinely
joint prerequisites occur; the rest is handled in polynomial time.

\section{The native loop matrix, worked explicitly}\label{sec:example}
The maintained source experiment gives five rows on the retained one-handle
labels $(3,5,6,7,8)$ \cite{repo-results}:
\begin{equation}\label{eq:loop-matrix}
N=\begin{pmatrix}
0&0&1&0&0\\
1&1&0&0&0\\
0&1&0&1&1\\
1&0&1&1&1\\
0&1&1&1&0
\end{pmatrix}.
\end{equation}
There is one interior dual vertex. Rows $0,1,2,3$ join it to the boundary root;
row~4 is its loop. Thus the four singleton sets $\{0\},\ldots,\{3\}$ are
exactly the dual spanning trees.

The default choice deletes row~0. The remaining column totals are
$(2,3,2,3,2)$, so all remaining rows form a stopping set. By
Lemma~\ref{lem:stopping}, no change of positive-pivot order, including the use
of fill, can rescue that particular tree choice.

The delivered general solver, with its incumbent shortcut disabled, finds the
tree $\{1\}$ after 97 order/partition branches. Its positive certificate uses
\begin{equation}\label{eq:loop-pivots}
 (r,h)=(3,3),\ (2,8),\ (4,5),\ (0,6).
\end{equation}
Here $h$ is the original source column label, not its zero-based position in
\eqref{eq:loop-matrix}. To check the schedule, assign the tree row time $-1$ and
the pivot rows their times $0,1,2,3$:
\begin{center}
\begin{tabular}{rlll}
\toprule
Time & Pivot & Other source rows in its column & Their deletion times\\
\midrule
0 & $(3,3)$ & $1$ & $-1$\\
1 & $(2,8)$ & $3$ & $0$\\
2 & $(4,5)$ & $1,2$ & $-1,1$\\
3 & $(0,6)$ & $3,4$ & $0,2$\\
\bottomrule
\end{tabular}
\end{center}
Every pivot entry is one and every other incident row is earlier. One column,
label~7, remains. This is a source-derived matrix-level recovery, not a new run
of the original triangulation factory or Regina oracle.

The maintained experiment records another valid schedule for the same deleted
row, using $(4,7)$ in the third step instead of $(4,5)$. That schedule leaves
column~5 rather than column~7. This difference is not cosmetic when passive
curves are present: a different surviving cocore can change a later private-port
witness. Our native integration seam proposes the dual forest and then sends the
source through its existing downstream checker rather than pretending that the
local matrix certificate alone proves the whole cut.

\section{Certificates, resource semantics, and integration}\label{sec:implementation}
\subsection{Positive packets}
The local schema \code{dual-forest-positive-v1} contains a tree, its root-to-leaf
dual cancellation order, and the one/two pivot pairs. The independent module
\code{dualforest/verify.py} imports neither the producer, the dependency builder,
nor the matroid optimizer. It checks the tree graph and the endpoint multiplicity
of each dual move, then checks source timestamps for all pivots.

A row deleted by the tree has time $-1$. A pivot row has its listed time; every
other source row meeting its pivot column must have smaller time. The checker
also requires distinct pivot columns, exact row coverage, unit pivot entries,
and valid labels. This costs one incidence scan plus graph and label bookkeeping.
A malformed or incomplete proof is rejected. No topological verdict is emitted
by these routines.

\subsection{Optimality and negative packets}
For pair-column instances, \code{pair-defect-v1} includes the optimal tree,
additional seed rows, a finishing schedule, the claimed deficit, and obstruction
colors. The independent checker rebuilds the low dependency relation from the
source and computes its strongly connected blocks using a separate, simple
reachability construction. It verifies the upper and lower bounds in
Theorem~\ref{thm:deficit}. This verifier is polynomial, but its deliberately
simple repeated traversals are not claimed to be linear.

For a general negative answer, \code{dual-forest-negative-v1} lists one rank
obstruction per canonical branch. The checker regenerates all exceptional
orders and all restricted-growth set partitions. Missing, extra, or invalid
obstructions cause rejection. It performs no matroid-intersection search.
Verification of a complete negative may be as expensive in the parameter as the
search, which is consistent with the hardness boundary.

\subsection{Unknown is distinct from no}
The solver accepts a nonnegative \code{max_branches} cap. If it is exhausted
before a positive branch or exhaustive negative completion, the answer is
\code{unknown}, with no negative certificate. The cap counts intersection
subproblems, not seconds or total instructions; input reading and dependency
preparation have their own cost. A callback is polled at search and augmentation
boundaries; its exceptions propagate, rather than becoming negative answers.

Python integers retain exactness. The source factory rejects Boolean values
masquerading as integers, negative incidences, duplicate serialized columns,
nonexistent endpoints, disconnected dual graphs, and inconsistent row counts.
A test reads an incidence $2^{16384}$, saturates it to two, and checks the expected
answer. This checks the kernel's binary coefficient handling, not the size or
cost of an entire normal-surface certificate with that multiplicity.

\subsection{The native seam and what has not been executed}
The record-level adapter maps native \code{None} endpoints to the root and each
native three-handle ID to a distinct interior vertex. It preserves repeated
endpoints and the original two-handle row IDs. Source one-handle column labels
are preserved, including gaps. Four unit tests exercise the translation and
malformed-record rejection.

The optional function \code{propose_geometry} in
\code{integration/native_adapter.py} imports the pinned native producer,
\code{_forest}, \code{_occurrences}, and
\code{replay_core_three_handles}. It obtains the current endpoint records,
checks them through the existing independent native source replayer, forms the
fixed-primal-forest matrix, and searches for a useful tree. On local success it
emits the existing native \code{core-three-handle-reduction-v1} shape and calls
the native replayer again on that proposed forest. The returned geometry must
still pass all normal core, port, boundary-pattern, and whole-cut verification.

\textbf{Execution boundary:} record translation was tested here; the native
\code{propose_geometry} call was not. A full repository checkout and its native
geometry dependencies were not available in the execution environment. No claim
is made that the complete native test suite was rerun or that a production
recognizer was modified. The archive is an integration candidate with an
explicit source-replay gate, not a silently installed replacement.

\paragraph{Recommended rollout.}
Retain an already successful native route unchanged. Use the pair-column
selector when its hypothesis holds. On a small number of dual choices, bounded
literal trials may be cheaper than the general FPT kernel; the measurements below
make that clear. Otherwise use the exceptional-boundary selector with an
explicit budget and an exact fallback or inconclusive return. Record $b$, the
number of source components, clone counts, actual branches, proof size, and
complete-query timing. Only source replay and whole-query comparisons can
justify a production default change.

\section{Reproducible validation and measured tradeoffs}\label{sec:experiments}
\subsection{Executed checks}
The package is standard-library Python, executed with Python~3.13.5. The 54 unit
tests include directional incidence cases, loop obstructions, empty inputs,
shared clone colors, an exchange requiring reassignment, exact deficits,
source-timestamp mutation checks, missing negative branches, large binary
entries, resource interruption, and native record translation.

The separate audit does not use the selector to supply its expected answers.
Its small-instance oracle literally enumerates dual spanning trees and scans
matrix column totals. The independent-transversal oracle enumerates color
assignments and checks graphical independence directly. Deficit checks enumerate
both trees and extra deletion subsets. The positive-fill oracle recursively tries
all singleton positive pivots with the update \eqref{eq:positive-update}.

\begin{table}[htbp]\centering\small
\begin{tabular}{p{3.7in}rr}\toprule
Audit & Inputs & Mismatches\\\midrule
All three-row, two-column ternary matrices on all connected two-vertex dual multigraphs & 13,851 & 0\\
All ternary $3\times3$ matrices on the loop graph & 19,683 & 0\\
Mixed random matrices, $m\le10$, $b\le4$ & 2,000 & 0\\
Unrestricted random $5\times5$ matrices & 150 & 0\\
Matroid solver versus assignment enumeration & 1,000 & 0\\
Pair-column optimal deficits & 500 & 0\\
Peeling versus all positive-pivot orders & 19,683 & 0\\
\bottomrule\end{tabular}
\caption{Executed finite audits. The pair-column stage also checks all 22,830 seed subsets of its 500 inputs. Counts in different stages are different comparison tasks, not a single disjoint population.}
\end{table}


All 35,684 inputs in the first four rows receive exact answers, with zero
mismatches and no branch-capped unknowns. Every positive certificate and every
available negative certificate in that audit is replayed independently. The
largest recorded branch count there is 3,220. These counts do not prove the
asymptotic theorems; they audit the implementation against their finite
consequences. Randomized tests are seeded and fully reproducible.

A separate hardness audit generates 400 CNF formulas and checks all 7,988 of
their assignments against the corresponding tree deletions. Forty additional
small formulas are compared with the exact general selector, and its positive
or negative certificates are replayed. There are zero mismatches. These are
checks of the explicit reduction, not experimental evidence of geometric
realizability.

\subsection{A family with exponentially many dual trees}
For $q$ arms, take two parallel edges $a_i,b_i$ from the root to vertex $i$.
Let row $a_i$ have a private singleton unit column and let row $b_i$ be zero.
The graph has $2^q$ spanning trees. The sole successful one chooses every $b_i$,
so that the anchored rows $a_i$ remain for peeling. The pair-column selector
identifies the singleton mandatory blocks $\{b_i\}$ and selects their
independent representatives without enumerating trees.

With the enumeration order used in the benchmark, that successful tree is last.
Swapping the row roles gives an easy-first control with a successful first tree.
Adding a zero root-loop row makes every tree fail and gives deletion deficit one.
These three families distinguish an asymptotic enumeration obstruction from
mere lucky or unlucky first choices. They are abstract matrices, not asserted
normal-cut families.

\subsection{Paired matrix-stage measurements}
Both arms receive the same already parsed instance. One runs the pair-column
selector; the other enumerates only \emph{actual} spanning trees of this star
family, not non-tree edge subsets, and applies literal peeling. Positive proof
generation and replay are included. The pair solver additionally replays its
optimality packet on negatives; the exhaustive negative baseline proves its
answer by finishing its enumeration. Source construction and all native topology
are excluded from both arms.

Each comparison has five randomly ordered pairs after an untimed warmup and
three identical-selector/identical-selector control pairs. The repository matrix
uses seven measurement pairs. All raw nanosecond-based samples, orders, and
controls are retained in \code{results/benchmark.json}. A ratio above one favors
the selector; the ratio is the median of paired ratios, not the quotient of
medians.

\begin{table}[htbp]\centering\small
\begin{tabular}{lrrrr}\toprule
Input & Selector (ms) & Trees (ms) & Ratio & A/A\\\midrule
Unique last, $q=1$ & 0.039 & 0.016 & 0.398 & 0.967\\
Unique last, $q=4$ & 0.068 & 0.098 & 1.375 & 0.966\\
Unique last, $q=8$ & 0.135 & 3.697 & 28.222 & 0.941\\
Unique last, $q=12$ & 0.250 & 136.882 & 553.569 & 1.057\\
Easy first, $q=12$ & 0.181 & 0.074 & 0.403 & 0.736\\
Obstructed, $q=12$ & 0.359 & 134.829 & 381.331 & 0.939\\
Repository loop matrix & 1.062 & 0.032 & 0.031 & 1.029\\
\bottomrule\end{tabular}
\caption{Matrix-stage medians and paired ratios, not native recognition timings. The loop row uses the general FPT kernel; the star rows use the polynomial pair-column kernel. A/A compares two identical selector executions.}
\end{table}


On the unique-last $q=12$ family the measured paired ratio is approximately
\UniqueLastRatio{}, and the obstructed family gives approximately \ObstructedRatio{}. These large factors
come from avoiding 4,096 enumerated trees in a deliberately specified structural
family. They are \emph{not} speedup claims for a knot recognizer.

The negative results matter equally. The easy-first $q=12$ baseline is faster,
with paired ratio about $\EasyFirstRatio$. For the five-row repository matrix the general
FPT selector has ratio about $\LoopRatio$: literal tree enumeration is roughly
\LoopSlowdown{} times faster on this tiny example. The selector solves a scalable structural
problem, but its first implementation should not displace inexpensive bounded
trials. The matrix example demonstrates recovered local coverage, not a
wall-clock improvement in native recognition.

\section{What this proves toward quasi-polynomial recognition}\label{sec:global}
Let a correct surrounding algorithm make $Q(n)$ authenticated core-selection
queries on an $n$-crossing input. For query $i$, let $\Lambda_i$ be its actual
encoded input length and $b_i$ its exceptional row boundary. Let $C_{\rm other}(n)$
charge every other task, including useful surface discovery, normal-coordinate
and pattern transport, source authentication, geometric rebuilding, reset work,
and any fallback outside this selector.

\begin{theorem}[Explicit composition budget]\label{thm:composition}
If the surrounding algorithm needs exactly the local decision in
Definition~\ref{def:problem}, and its correctness is maintained by the source
and downstream checks described above, then replacing its forest selection by
Theorem~\ref{thm:fpt} gives total cost at most
\begin{equation}\label{eq:global-budget}
 C_{\rm other}(n)+\sum_{i=1}^{Q(n)}
       \poly(\Lambda_i)\,2^{O(b_i\log(b_i+1))}.
\end{equation}
In particular, the target $n^{O(\log n)}$ follows from the sufficient conditions
\begin{align}
 C_{\rm other}(n),\ Q(n)&\le 2^{O((\log n)^2)},\label{eq:cond-other}\\
 \Lambda_i&\le n^{O(1)},\label{eq:cond-size}\\
 b_i\log(b_i+1)&=O((\log n)^2)\quad\text{uniformly in }i.\label{eq:cond-b}
\end{align}
For example, $b_i=O((\log n)^2/\log\log n)$ suffices for large $n$.
\end{theorem}
\begin{proof}
Each query cost is Theorem~\ref{thm:fpt}; all omitted work is charged to
$C_{\rm other}$. Under the stated conditions each summand and the number of
summands are bounded by $2^{O((\log n)^2)}$. Their product has the same form.
The example for $b_i$ follows from $\log(b_i+1)=O(\log\log n)$.
\end{proof}

This theorem is an accounting implication, not a proof of its hypotheses.
No universal bound on the realized $b_i$, source sizes, or reset cost is established
here. In particular, a local successful tree may leave a different passive disk
system and alter later costs. Neither the pair-column theorem nor the hardness
construction resolves that geometric issue.

The iteration bound $L(g+1)^L$ in Section~9 of Lackenby's preprint
\cite{lackenby2026} highlights why local search cost and hierarchy rebuild cost
must be kept separate. The new selector removes spanning-tree enumeration in its
polynomial case and confines its own nonpolynomial work to an explicit higher-
support row boundary. It does not by itself control $L$, $g$, or the costs of the
other operations.

\section{Further research questions and concrete next steps}\label{sec:questions}
\begin{question}[Measure the source boundary, not only tetrahedra]
What are the distributions of $b$, low-graph source components, and graphical
rank deficits on the maintained authenticated cut corpus, including historical
source cores and the later loop/interior-vertex fixtures?
\end{question}
The next experiment should export the raw counts and dual endpoints from the
native source checker, then record both kernel time and completed whole-cut time.
The historical native corpus was not rerun in this package. A useful outcome would identify whether difficult cases have a
small exceptional boundary or whether a different representation is needed.

\begin{question}[Choose the primal forest to reduce the exceptional boundary]
Can one optimize or approximate $b$ over authenticated zero/one-handle forests,
while keeping the dual endpoints and the later port data source-bound?
\end{question}
The current theory deliberately freezes that forest. A primal change can remove
columns and alter which columns have high support; treating those changes as
free preprocessing would conceal a second search. A useful theorem would combine
primal selection with a controlled, explicitly charged change in the row boundary.

\begin{question}[Compress exceptional orders]
Can the $b!\Bell(b)$ factor be replaced by $2^{O(b)}$, perhaps through a
representative-family or subset dynamic program that preserves shared witnesses?
\end{question}
The obstruction is not merely enumerating permutations: multiple ancestor
requirements may share one edge, and the chosen edges must remain jointly
independent with many mandatory blocks. A correct compression must preserve those
simultaneous constraints, not only whether each demand has some candidate.

\begin{question}[Parameterize high targets rather than all high-support rows]
Can one replace $b$ by the number of rows which carry a unit entry in a
higher-support column, plus a smaller interface parameter for its nonunit rows?
\end{question}
A nonunit row can never be the high-column pivot, but its deletion is still a
prerequisite. The present proof includes it in $B$ so that the prefix test is
complete. Omitting it without a substitute would invalidate the order theorem.
A succinct representation of these prerequisites could be especially valuable
for large sparse geometric interfaces.

\begin{question}[What restrictions does genuine geometry impose?]
Which matrix-and-dual-graph pairs are realizable by the authenticated source core
construction, and do those restrictions exclude the satisfiability gadgets of
Theorem~\ref{thm:hardness}?
\end{question}
A realization theorem or an obstruction theorem would both be useful. Until
one is proved, the abstract NP-completeness result must not be promoted to a
hardness statement about normal cuts. Conversely, dismissing it without proving
an appropriate source restriction would not justify a general polynomial
selector.

\begin{question}[Turn deletion deficits into legal geometric progress]
Given an optimal deficit certificate, can the indicated unhit source blocks be
resolved by a bounded-cost collection of authenticated disk changes or local
moves, with a provable decrease in a global progress measure?
\end{question}
The additional rows $Z$ are not legal deletions merely because the optimizer
names them. A useful bridge would associate a geometric operation and an explicit
cost or obstruction with each rank-deficient group, rather than using the deficit
as an informal heuristic.

\begin{question}[Optimize downstream witnesses, not just raw completion]
Among successful dual forests, can one select a forest whose resulting cocore
system preserves enough private ports or pattern witnesses to minimize later
rebuilding?
\end{question}
The loop example already has different valid pivot schedules leaving different
columns. A weighted representative objective may help when costs are additive,
but passive-witness compatibility need not be additive or matroidal. A sound
state description and source-bound transport theorem are required before a
weighted optimization can claim a hierarchy benefit.

\begin{question}[Count-aware or bounded-trial portfolios]
Can a cheap bound on the number of dual trees choose reliably between bounded
literal enumeration and the FPT kernel, while preserving a clean worst-case
budget?
\end{question}
The benchmark suggests this immediately: the small native example favors
literal trials, while the star families favor matroid selection. A fixed number
of verified trials is easy to charge before the exact selector. More ambitious
selection could use an exact tree count and an output-sensitive tree enumerator;
those components would need their own bit and certificate accounting.

\begin{question}[Stream and compress negative branch proofs]
Can identical rank obstructions across different orders and witness partitions
be shared without requiring the checker to trust a producer's branch omissions?
\end{question}
The current negative checker deliberately regenerates every branch. A compressed
proof would need a verifiable coverage argument for the omitted order regions.
This may lead to practical reductions well before a sharper theoretical FPT bound.

\begin{question}[Establish a source-aware quasi-polynomial envelope]
Can the hierarchy construction be arranged so that the actual source sizes,
exceptional boundaries, and number of rebuilds satisfy
\eqref{eq:cond-other}--\eqref{eq:cond-b} simultaneously?
\end{question}
The simultaneous qualification matters. A representation reducing $b$ may enlarge
the source or disrupt a previously short hierarchy. A theorem controlling only
one quantity does not close the recognition bound. The most consequential next
step is a geometric normal form that preserves source authentication while
controlling all three quantities over the entire run.

\section{Conclusion}
A complete dual cancellation forest is not necessarily a useful one. For the
pair-column regime, the correct missing condition is an independent transversal
of unanchored source components, and this gives both a polynomial selector and
an exact deletion-deficit certificate. For general cores, the higher-support row
boundary is enough to localize the nonpolynomial search: exceptional orders and
shared-witness partitions reduce every branch to the same checkable matroid
primitive. The support-three hardness construction shows why a general abstract
polynomial claim would require more structure.

The delivered kernel recovers the documented matrix-level loop miss and passes
its finite audits, while the benchmark honestly retains cases where direct tree
trials are faster. Its immediate role is an opt-in, source-replayed selection
component and a measurable structural parameter for the maintained pipeline.
The route to unrestricted quasi-polynomial recognition still needs a theorem
controlling the source geometry, exceptional boundaries, and hierarchy resets
jointly.

\appendix
\section{Algorithm pseudocode and certificate interfaces}
\subsection{The pair-column selector}
\begin{lstlisting}[language=Python]
G, anchors = low_column_dependencies(N)
blocks = source_SCCs(G) without anchors
I, rank_witness = maximum_independent_transversal(D, blocks)
T = extend_to_spanning_tree(D, I)
Z = one row from each block not met by T
pivots = peel(N, initially_deleted=T | Z)
return T, Z, pivots, rank_witness
\end{lstlisting}
When $Z$ is empty this is a positive selection proof. Otherwise the rank witness
and complete schedule after $T\cup Z$ prove the optimum deletion deficit.

\subsection{The general selector}
\begin{lstlisting}[language=Python]
G, anchors, high_columns, B = classify_usable_columns(N)
U = source_SCCs(G) containing neither an anchor nor a row of B
R0 = reach(G, anchors | union(U))
for order in permutations(B):
    Q = nonautomatic_ancestor_requirements(order, R0, G, high_columns)
    for partition in set_partitions(indices(Q)):
        groups = U + [intersection(Q[i] for i in block)
                      for block in partition]
        I, obstruction = maximum_independent_transversal(D, groups)
        if I represents every group:
            T = extend_to_spanning_tree(D, I)
            return YES, T, peel(N, T)
        record(obstruction)
return NO, complete_branch_obstruction_record
\end{lstlisting}
A cap checked before each intersection call replaces the last line by
\code{UNKNOWN} when exhausted. No bound on source preparation is implied by that
branch cap.

\subsection{Reproduction commands}
From the extracted package root:
\begin{lstlisting}[language=bash]
python -m unittest discover -s tests -v
python -m experiments.validate
python -m experiments.sat_audit
python -m experiments.benchmark
python -m dualforest examples/repository_loop.json --max-branches 10000
sh article/build.sh
\end{lstlisting}
The example command consumes a normalized matrix/dual-graph input, not a knot
diagram. It cannot certify unknotness by itself. The JSON packet explicitly
records its local scope.

\section{Source and artifact provenance}
The primary source pin is \pin. The mathematical source readings were the full
maintained peeling, dual-three-handle, and dual-result sections; the native
three-handle producer and independent checker were also read in full. The native
\code{core_handlebody.py} reading covered its input contract, primal forest,
occurrence extraction, and the beginning of its capped planner. An indexed
search confirmed the current column-peeling entry point, but the whole native
module was not imported or executed.

The incoming directory was inventoried. In particular,
\code{ProveIt_Dual_Forest_Normalization_2026-10-12.zip} and
\code{ProveIt_Dual_Forests_Homological_Progress_2026-10-12.zip}
were present, but their internal manuscripts and code were not read. The present
report must be reconciled with those deliveries during intake before any claim
of repository-wide novelty. No incoming files or maintained repository files were
changed by this session.

The included examples, source-neutral implementation, synthetic reduction,
finite-oracle code, recorded outputs, and this article were generated for the
present package. The five-row matrix is explicitly attributed to the maintained
source result. External research papers and Lackenby's slide deck are cited but
not redistributed. \code{README.md} gives the file map and integration
instructions; \code{CLAIMS.md} separates proven, tested, inherited, conditional,
and unexecuted items.

\begin{thebibliography}{99}
\bibitem{repo-peeling}
V.~Reshetnikov and contributors.
\emph{Complete rectangular cancellation by sparse column peeling}.
ProveIt, \code{Topology/UnknotRecognition/synthesis/column_cancellation.tex},
commit \pin, 2026.
\url{https://github.com/VladimirReshetnikov/ProveIt}.

\bibitem{repo-dual}
V.~Reshetnikov and contributors.
\emph{Removing source core three-handles relative to the bundle ports}.
ProveIt, \code{Topology/UnknotRecognition/synthesis/core_three_handles.tex},
same pinned commit, 2026.

\bibitem{repo-results}
V.~Reshetnikov and contributors.
\emph{Source tests, dual-forest dependence, and three-handle results}.
ProveIt, \code{Topology/UnknotRecognition/synthesis/three_handles_results.tex},
same pinned commit, 2026. The five-row loop matrix and original alternative
certificates are recorded here.

\bibitem{repo-native}
V.~Reshetnikov and contributors.
\emph{Source-authenticated removal of core three-handles by a dual boundary forest}.
ProveIt, \code{Topology/UnknotRecognition/fast/fastunknot/core_three_handle_reduction.py},
same pinned commit, 2026.

\bibitem{repo-checker}
V.~Reshetnikov and contributors.
\emph{Independent replay of the source core's dual two/three cancellation forest}.
ProveIt, \code{Topology/UnknotRecognition/fast/fastunknot/core_three_handle_reduction_verify.py},
same pinned commit, 2026.

\bibitem{lackenby-talk}
M.~Lackenby.
\emph{Quasi-polynomial unknot-recognition lecture slides}.
Oxford-hosted slide deck, inspected 12 October 2026; displayed main theorem at
PDF page~13.
\url{https://people.maths.ox.ac.uk/lackenby/quasipolynomial-talk.pdf}.

\bibitem{lackenby2026}
M.~Lackenby.
\emph{Incompressible surfaces, hierarchies and unknot recognition}.
arXiv:2607.23350v1, 2026, especially Section~9 and Proposition~9.1.
\url{https://arxiv.org/html/2607.23350v1}.

\bibitem{matroid2024}
M.~B\'ar\'asz, K.~B\'erczi, T.~Kir\'aly, Y.~Yamaguchi, and Y.~Yokoi.
\emph{Matroid Intersection under Minimum Rank Oracle}.
arXiv:2407.03229v1; EGRES Technical Report TR-2024-13, 2024.
The version is specified because later versions have changed author metadata.
\url{https://arxiv.org/html/2407.03229v1}.

\bibitem{cook1971}
S.~A.~Cook.
The complexity of theorem-proving procedures.
In \emph{Proceedings of the Third Annual ACM Symposium on Theory of Computing},
pages 151--158, 1971.
\url{https://doi.org/10.1145/800157.805047}.
\end{thebibliography}
\end{document}
